% Volume II: Causal Explosion
% Part IV. Time, Delay, and the Maturation of Evidence
% Chapter 20: Discovery Curves
% Spine: proof-spines/ch020-spine.md

\chapter{Discovery Curves}
\label{ch:ch020}

This chapter models \textbf{discovery curves} by letting \(F(t)\) denote the fraction of ultimately relevant evidence recognized by time \(t\). Different inquiries produce different shapes: exponential, logistic, power-law, and punctuated curves all describe ways in which recognition can accumulate. Sometimes a late archival find or forensic breakthrough creates a discontinuity,
\[
F(t^\ast+)-F(t^\ast-)\gg 0,
\]
\ToyModel\ This jump condition sketches a late-breakthrough discovery curve without claiming a universal fit.
showing that evidentiary maturation is not always smooth.

The institutional problem is how to behave when \(F(t)\) remains far from saturation. Slow confirmation is not always weakness, and rapid early uptake is not always vindication. Resource allocation, investigative persistence, and reopening rules should be tuned to the expected shape of the curve in the relevant domain, because the tempo of learning is itself part of the evidence.
